\subsection{Проверка устойчивости с помощью годографа Михайлова}

Для проверки устойчивости системы используем критерий
Михайлова.

Характеристический полином системы имеет вид

\begin{equation}
A(p)=a_0p^6+a_1p^5+a_2p^4+a_3p^3+a_4p^2+a_5p+a_6.
\end{equation}

Выполним замену
\[
p=i\omega.
\]

Тогда характеристический полином примет вид

\begin{equation}
A(i\omega)=\operatorname{Re}(\omega)+
i\operatorname{Im}(\omega),
\end{equation}

где вещественная и мнимая части определяются выражениями

\begin{equation}
\operatorname{Re}(\omega)=
a_6-a_4\omega^2+a_2\omega^4-a_0\omega^6,
\end{equation}

\begin{equation}
\operatorname{Im}(\omega)=
a_5\omega-a_3\omega^3+a_1\omega^5.
\end{equation}

Согласно критерию Михайлова, для устойчивой системы
годограф характеристического полинома должен при
\[
\omega:0\to+\infty
\]
совершить поворот на угол
\[
\frac{n\pi}{2},
\]
где \(n\) — порядок характеристического полинома.

Так как порядок полинома равен шести,
годограф должен совершить поворот на угол

\[
3\pi.
\]

Для проверки критерия выберем точку области устойчивости
\[
k_{\text{п}}=10,
\qquad
k_{\text{и}}=15.
\]

Тогда коэффициенты характеристического полинома принимают
конкретные значения, а вещественная и мнимая части
годографа Михайлова имеют вид:

\begin{equation}
\begin{aligned}
\operatorname{Re}(\omega) =\;&
-0.015\omega^6
+ 10542\omega^4
- 1.05154 \cdot 10^7 \omega^2
+ 5.25 \cdot 10^7,
\\
\operatorname{Im}(\omega) =\;&
5.072\omega^5
- 48100\omega^3
+ 3.521 \cdot 10^7 \omega.
\end{aligned}
\end{equation}

Для проверки устойчивости системы был построен годограф Михайлова для характеристического полинома шестого порядка. Построение проводилось при различных диапазонах изменения частоты $\omega$, поскольку при больших значениях частоты масштаб графика существенно изменяется.

На рисунке~\ref{fig:pdf/mihailov_1} представлен годограф при диапазоне
\[
\omega \in [0.01;1].
\]
В данном диапазоне исследуется поведение годографа вблизи начала координат.

\screenshot[0.75]{pdf/mihailov_1}{Годограф Михайлова при $\omega \in [0.01;1]$}

На рисунке~\ref{fig:pdf/mihailov_10} показан годограф при диапазоне
\[
\omega \in [0.01;10].
\]
При увеличении диапазона частот становится заметным дальнейшее движение траектории в комплексной плоскости.

\screenshot[0.75]{pdf/mihailov_10}{Годограф Михайлова при $\omega \in [0.01;10]$}

На рисунке~\ref{fig:pdf/mihailov_5000} представлен годограф при диапазоне
\[
\omega \in [0.01;5000].
\]
В этом случае хорошо видно, что годограф продолжает поворот против часовой стрелки и уходит в левую полуплоскость.

\screenshot[0.75]{pdf/mihailov_5000}{Годограф Михайлова при $\omega \in [0.01;5000]$}

Полученные результаты показывают, что годограф Михайлова совершает поворот, соответствующий порядку характеристического полинома. Следовательно, исследуемая система удовлетворяет критерию Михайлова и является устойчивой.
